% Tabla G1: Lecturas de calibre de las réplicas.
% Etiqueta original: tab:lecturas-calibre.
% Fuente: anexos/anexos.tex, lineas 174-192.
% Fragmento extraido automaticamente; requiere el preambulo de main.tex.

\begin{table}[htbp]
\centering\footnotesize
\caption{Lecturas de calibre de las réplicas (mm).}
\label{tab:lecturas-calibre}
\begin{tabular}{@{}lccccc@{\hspace{1.2em}}ccccc@{}}
\toprule
 & \multicolumn{5}{c}{Longitud total (LT)} & \multicolumn{5}{c}{Profundidad corporal (PC)} \\
\cmidrule(lr){2-6}\cmidrule(l){7-11}
Réplica & 1 & 2 & 3 & Media & Desv. & 1 & 2 & 3 & Media & Desv. \\
\midrule
R1 & 160 & 157 & 159 & 158,7 & 1,5 & 45 & 45 & 45 & 45,0 & 0,0 \\
R2 & 131 & 136 & 136 & 134,3 & 2,9 & 28 & 27 & 27 & 27,3 & 0,6 \\
R3 & 91  & 92  & 92  & 91,7  & 0,6 & 41 & 40 & 40 & 40,3 & 0,6 \\
R4 & 141 & 146 & 141 & 142,7 & 2,9 & 43 & 42 & 43 & 42,7 & 0,6 \\
R5 & 137 & 138 & 137 & 137,3 & 0,6 & 25 & 26 & 26 & 25,7 & 0,6 \\
R6 & 210 & 209 & 210 & 209,7 & 0,6 & 35 & 36 & 35 & 35,3 & 0,6 \\
\bottomrule
\end{tabular}
\end{table}
